\lecture{1}{Introduction to Machine Learning with Graphs}{Tuesday, September 22, 2026}\label{ch:lec01}

\textbf{Outline.}  why graphs (graphs as a language, databases,
why graph DL is hard) $\to$ choice of graph representation $\to$ applications: node-,
link- and graph-level tasks.


% =====================================================================
\section{Why graphs?}\label{sec:lec01-why-graphs}

\subsection{Graphs: a general language}\label{sec:lec01-language}
\slides{22--25}

\begin{keyidea}
  Graphs are a general language for describing and analyzing \textbf{entities with
  relations/interactions}.
\end{keyidea}

\textbf{Many types of data are graphs:}
\begin{itemize}
  \item event graphs, computer networks, disease pathways, food webs, particle networks,
    underground networks;
  \item social, economic, communication and citation networks, the Internet, networks of
    neurons;
  \item knowledge graphs (e.g.\ Leonard Nimoy $\xrightarrow{\text{played}}$ Spock
    $\xrightarrow{\text{characterIn}}$ Star Trek $\xrightarrow{\text{genre}}$ Science
    Fiction), regulatory networks, scene graphs, code graphs, molecules, 3D shapes.
\end{itemize}

\subsection{Databases are graphs: Relational Deep Learning}\label{sec:lec01-rdl}
\slides{26--28}

\begin{itemize}
  \item Relational databases = tables linked by keys: \textbf{commerce} (customers, orders,
    orders\_items, items), \textbf{finance} (loan, order, trans, account, disp, card,
    client, district), \textbf{clinical trial} (patients, admissions, noteevents,
    procedures\_icd, diagnoses\_icd).
  \item \term{Relational Deep Learning} (RDL): tables Users / Sales (ID, Product\_ID,
    Session\_ID, User\_ID, Price) / Products $\Rightarrow$ graph with \textbf{rows as
    nodes} and \textbf{key links as edges} (\cref{fig:lec01-rdl}). Benchmark:
    \texttt{relbench.stanford.edu}.
\end{itemize}

\begin{figure}[htbp]
  \centering
  \begin{tikzpicture}[scale=0.85, transform shape,
      tbl/.style={draw, rounded corners=2pt, thick, fill=black!4, align=left,
                  font=\scriptsize, inner sep=3pt},
      usr/.style={circle, draw=black!70, thick, fill=Orchid!35, minimum size=5mm,
                  inner sep=0pt, font=\scriptsize},
      sale/.style={rectangle, rounded corners=1.5pt, draw=black!70, thick,
                   fill=Goldenrod!35, inner sep=1.5pt, font=\tiny},
      prod/.style={circle, draw=black!70, thick, fill=SkyBlue!45, minimum size=5mm,
                   inner sep=0pt}]
    \node[tbl] (ts) at (2.1,1.4) {\textbf{Sales}\\ ID\\ Product\_ID\\ Session\_ID\\ User\_ID\\ Price};
    \node[tbl] (tu) at (-0.6,1.4) {\textbf{Users}\\ ID};
    \node[tbl] (tp) at (5.0,1.4) {\textbf{Products}\\ Product\_ID};
    \draw[arrow, black!60] (ts.west) -- (tu.east);
    \draw[arrow, black!60] (ts.east) -- (tp.west);
    \draw[-{Stealth[length=3mm]}, line width=1.6pt, black!40] (6.2,1.4) -- (6.9,1.4);
    \foreach \i/\y in {1/3.0, 2/2.0, 3/1.0, 4/0.0} \node[usr] (u\i) at (7.4,\y) {$u_\i$};
    \foreach \i/\y/\p in {1/3.3/20, 2/2.75/80, 3/2.2/15, 4/1.65/20, 5/1.1/80, 6/0.55/80, 7/0.0/15}
      \node[sale] (s\i) at (9.1,\y) {\$\p};
    \node[prod, label={[font=\scriptsize]right:shirt}] (shirt) at (10.8,2.75) {};
    \node[prod, label={[font=\scriptsize]right:shoe}]  (shoe)  at (10.8,1.55) {};
    \node[prod, label={[font=\scriptsize]right:sunglasses}] (sun) at (10.8,0.35) {};
    \draw[edge] (u1)--(s1) (u1)--(s2) (u2)--(s3) (u3)--(s4) (u3)--(s5) (u4)--(s6) (u4)--(s7);
    \draw[edge] (s1)--(shirt) (s4)--(shirt) (s2)--(shoe) (s5)--(shoe) (s6)--(shoe)
                (s3)--(sun) (s7)--(sun);
    \node[font=\scriptsize] at (7.4,-0.6) {Users};
    \node[font=\scriptsize] at (9.1,-0.6) {Sales};
    \node[font=\scriptsize] at (10.8,-0.6) {Products};
  \end{tikzpicture}
  \caption{RDL: schema $\to$ graph of rows (slide 27; prices from slides 46/68). Reused
  for the churn example and the summary.}\label{fig:lec01-rdl}
\end{figure}

\begin{keyidea}
  Complex domains have a rich relational structure, which can be represented as a
  \textbf{relational graph}. \textbf{By explicitly modeling relationships we achieve
  better performance!}\\
  \textbf{Main question:} how do we take advantage of relational structure for better
  prediction?
\end{keyidea}

\subsection{Why is graph deep learning hard?}\label{sec:lec01-hard}
\slides{29--32}

\begin{itemize}
  \item \textbf{Modern deep learning toolbox is designed for simple sequences \& grids}:
    images = grids (CNN: patterns of local contrast $\to$ face features $\to$ faces);
    text/speech/audio = sequences (RNN/LSTM).
  \item \textbf{Networks are complex:}
    \begin{itemize}
      \item arbitrary size and complex topological structure (i.e.\ \textbf{no spatial
        locality} like grids);
      \item \textbf{no fixed node ordering} or reference point;
      \item often \textbf{dynamic} and have \textbf{multimodal features}.
    \end{itemize}
\end{itemize}

\begin{keyidea}[This course]
  How can we develop neural networks that are much more broadly applicable?
  \textbf{Graphs are the new frontier of deep learning.}
\end{keyidea}

% =====================================================================
\section{Choice of graph representation}\label{sec:lec01-representation}

\subsection{Graphs: a common language}\label{sec:lec01-common}
\slides{34}

\begin{itemize}
  \item The \textbf{same graph} describes different domains: actors linked by movies
    (Movie 1, 2, 3), people linked by relations (friend, brothers, co-worker), interacting
    proteins (1, 2, 5, 9), or abstract nodes $1$--$4$ with $|N|=4$, $|E|=4$
    (\cref{fig:lec01-common}). (Slide uses $N$ for the node set; from slide 35 on it is $V$.)
\end{itemize}

\begin{figure}[htbp]
  \centering
  \begin{tikzpicture}[scale=0.85, lab/.style={font=\scriptsize\itshape, midway}]
    \node[vtx, fill=Red!40, label=above:{\small Peter}]  (P) at (-0.4,1.3) {};
    \node[vtx, fill=Red!40, label=above:{\small Mary}]   (M) at (2.2,1.45) {};
    \node[vtx, fill=Red!40, label=right:{\small Tom}]    (T) at (4.0,0.8) {};
    \node[vtx, fill=Red!40, label=below:{\small Albert}] (A) at (1.2,0) {};
    \draw[edge] (P) -- node[lab, above]{friend} (M);
    \draw[edge] (P) -- node[lab, left]{brothers} (A);
    \draw[edge] (M) -- node[lab, left]{friend} (A);
    \draw[edge] (M) -- node[lab, below, sloped]{co-worker} (T);
    \node[vertex] (v1) at (6.0,1.3) {1};
    \node[vertex] (v2) at (8.2,1.45) {2};
    \node[vertex] (v4) at (9.8,0.9) {4};
    \node[vertex] (v3) at (7.4,0) {3};
    \draw[edge] (v1)--(v2) (v1)--(v3) (v2)--(v3) (v2)--(v4);
    \node[font=\small, align=left] at (9.6,-0.1) {$|N| = 4$\\ $|E| = 4$};
  \end{tikzpicture}
  \caption{One structure, many domains (slide 34): social network (left) and the abstract
  graph (right); the actor and protein networks are identical.}\label{fig:lec01-common}
\end{figure}

\subsection{Heterogeneous graphs}\label{sec:lec01-hetero}
\slides{35--36}

\begin{definition}[Heterogeneous graph]\label{def:lec01-hetero}
  A \term{heterogeneous graph} is defined as
  \begin{equation}\label{eq:lec01-hetero}
    G = (V, E, R, T)
  \end{equation}
  \begin{itemize}
    \item nodes with node types $v_i \in V$;
    \item edges with relation types $(v_i, r, v_j) \in E$;
    \item node type $T(v_i)$; relation type $r \in R$;
    \item nodes and edges have \textbf{attributes/features}.
  \end{itemize}
\end{definition}

\textbf{Many graphs are heterogeneous:}
\begin{itemize}
  \item \textbf{Biomedical knowledge graphs:} node types drug, disease, adverse event,
    protein, pathway; relations e.g.\ TREATS, CAUSES, IS\_A, ASSOC, HAS\_FUNC.
    Example node: \emph{Migraine}; edge: (fulvestrant, Treats, Breast Neoplasms); node type:
    \emph{Protein}; edge type (relation): \emph{Causes}.
  \item \textbf{Academic graphs} (\cref{fig:lec01-academic}): example node: \emph{ICML};
    edge: (GraphSAGE, NeurIPS); node type: \emph{Author}; edge type (relation):
    \emph{pubYear}.
\end{itemize}

\begin{figure}[htbp]
  \centering
  \begin{tikzpicture}[scale=0.85, transform shape,
      sch/.style={draw, thick, minimum width=14mm, minimum height=5.5mm, font=\scriptsize,
                  fill=white},
      rl/.style={font=\tiny, text=BrickRed}]
    \node[sch] (Pa) at (0,0) {Paper};
    \node[sch] (Co) at (-3.0,0) {Conference};
    \node[sch] (Ti) at (2.9,0) {Title};
    \node[sch] (Au) at (-1.2,-1.5) {Author};
    \node[sch] (Ye) at (1.2,-1.5) {Year};
    \draw[dedge] (Pa) -- node[rl, above]{pubWhere} (Co);
    \draw[dedge] (Pa) -- node[rl, above]{hasTitle} (Ti);
    \draw[dedge] (Pa) -- node[rl, left]{hasAuthor} (Au);
    \draw[dedge] (Pa) -- node[rl, right]{pubYear} (Ye);
    \draw[dedge] ([xshift=-2mm]Pa.north) to[out=110, in=70, looseness=8]
      node[rl, above, font=\scriptsize]{cite} ([xshift=2mm]Pa.north);
  \end{tikzpicture}
  \caption{Academic graph schema: node types and relation types (slide 36).}\label{fig:lec01-academic}
\end{figure}

\begin{caution}
  Slide 36 labels the relation ``hasTtile'' (typo for \emph{hasTitle}); its example edge
  (GraphSAGE, NeurIPS) omits the relation -- as a triple: (GraphSAGE, pubWhere, NeurIPS).
\end{caution}

\subsection{Choosing a proper representation}\label{sec:lec01-choosing}
\slides{37}

\begin{itemize}
  \item How to build a graph: \textbf{what are nodes? what are edges?}
  \item \textbf{Choice of the proper network representation of a given domain/problem
    determines our ability to use networks successfully:}
    \begin{itemize}
      \item in some cases there is a unique, unambiguous representation;
      \item in other cases the representation is by no means unique;
      \item the way you assign links will determine the nature of the question you can
        study.
    \end{itemize}
\end{itemize}

\subsection{Directed vs.\ undirected graphs}\label{sec:lec01-directed}
\slides{38}

\begin{itemize}
  \item \term{Undirected}: links undirected (\textbf{symmetrical, reciprocal}).
    \term{Directed}: links directed (\cref{fig:lec01-dir}).
  \item Other considerations: \textbf{weights}, \textbf{properties}, \textbf{types},
    \textbf{attributes}.
\end{itemize}

\begin{figure}[htbp]
  \centering
  \begin{tikzpicture}[scale=0.72,
      rv/.style={circle, fill=Red!70, draw=Red!80!black, inner sep=0pt, minimum size=2.8mm},
      gv/.style={circle, fill=Green!60!black, draw=Green!40!black, inner sep=0pt, minimum size=2.8mm},
      nl/.style={font=\scriptsize\bfseries}]
    % undirected
    \node[rv, label={[nl]left:A}]  (A) at (2.1,-1.1) {};
    \node[rv, label={[nl]left:L}]  (L) at (4.2,-0.8) {};
    \node[rv, label={[nl]right:M}] (M) at (6.55,-1.4) {};
    \node[rv, label={[nl]above:F}] (F) at (4.7,-2.0) {};
    \node[rv, label={[nl]below right:D}] (D) at (2.9,-2.5) {};
    \node[rv, label={[nl]below:B}] (B) at (1.4,-2.85) {};
    \node[rv, label={[nl]above:I}] (I) at (7.35,-2.6) {};
    \node[rv, label={[nl]below left:G}] (G) at (5.7,-3.05) {};
    \node[rv, label={[nl]below:C}] (C) at (2.65,-4.15) {};
    \node[rv, label={[nl]below:H}] (H) at (6.75,-3.75) {};
    \draw[edge] (A)--(D) (B)--(D) (D)--(F) (F)--(G) (G)--(H) (H)--(I) (L)--(M);
    \draw[edge] (C) to[bend left=25] (D);
    \draw[edge] (C) to[bend right=25] (D);
    \draw[edge] (B) to[out=160, in=210, looseness=14] (B);
    \node[font=\small\bfseries, text=Magenta] at (4.2,0.0) {Undirected};
    % directed
    \begin{scope}[xshift=7.0cm]
      \node[gv, label={[nl]above:B}]  (b) at (3.8,-1.05) {};
      \node[gv, label={[nl]above:C}]  (c) at (5.6,-1.85) {};
      \node[gv, label={[nl]above:D}]  (d) at (7.1,-1.25) {};
      \node[gv, label={[nl]below:A}]  (a) at (2.8,-3.15) {};
      \node[gv, label={[nl]left:G}]   (g) at (5.4,-3.2) {};
      \node[gv, label={[nl]above:E}]  (e) at (8.7,-3.0) {};
      \node[gv, label={[nl]right:F}]  (f) at (7.2,-4.3) {};
      \draw[dedge] (a) to[bend left=18] (b);
      \draw[dedge] (b) to[bend left=18] (a);
      \draw[dedge] (b)--(c);
      \draw[dedge] (a)--(c);
      \draw[dedge] (d)--(c);
      \draw[dedge] (e)--(d);
      \draw[dedge] (f)--(e);
      \draw[dedge] (g)--(f);
      \node[font=\small\bfseries, text=Magenta] at (5.8,0.0) {Directed};
    \end{scope}
  \end{tikzpicture}
  \caption{Undirected vs.\ directed example graphs (slide 38); the undirected one includes a
  self-loop at B and a double edge C--D.}\label{fig:lec01-dir}
\end{figure}

\subsection{Bipartite graphs}\label{sec:lec01-bipartite}
\slides{39}

\begin{definition}[Bipartite graph]\label{def:lec01-bipartite}
  A \term{bipartite graph} is a graph whose nodes can be divided into two disjoint sets $U$
  and $V$ such that every link connects a node in $U$ to one in $V$; that is, $U$ and $V$
  are \textbf{independent sets}.
\end{definition}

\begin{minipage}[c]{0.7\linewidth}
\begin{itemize}
  \item \textbf{Examples:} Authors-to-Papers (they authored), Actors-to-Movies (they
    appeared in), Users-to-Movies (they rated), Recipes-to-Ingredients (they contain).
  \item \term{``Folded'' networks}: author collaboration networks, movie co-rating
    networks. (Folding = keep one side, link two of its nodes if they share a neighbour on
    the other side.)
\end{itemize}
\end{minipage}\hfill
\begin{minipage}[c]{0.26\linewidth}
  \centering
  \begin{tikzpicture}[scale=0.6, transform shape,
      uv/.style={vertex, fill=Red!45}, vv/.style={vertex, fill=Green!45}]
    \foreach \n/\y in {A/3.2, B/2.4, C/1.6, D/0.8, E/0.0} \node[uv] (\n) at (0,\y) {\n};
    \foreach \n/\y in {1/2.9, 2/2.0, 3/1.1, 4/0.2} \node[vv] (v\n) at (2,\y) {};
    \draw[edge] (A)--(v1) (A)--(v2) (B)--(v1) (B)--(v3) (C)--(v3) (D)--(v2) (D)--(v4)
                (E)--(v1) (E)--(v4);
    \node at (0,-0.7) {\Large $U$}; \node at (2,-0.7) {\Large $V$};
  \end{tikzpicture}\\[-2pt]
  {\footnotesize Slide 39 example.}
\end{minipage}

% =====================================================================
\section{Applications of graph ML: types of tasks}\label{sec:lec01-tasks}
\slides{40--42}

\begin{itemize}
  \item Task levels: \textbf{node level}; \textbf{edge level}; \textbf{community
    (subgraph) level}; \textbf{graph level} -- graph-level prediction, graph generation.
  \item Review: \textbf{node-level}, \textbf{link-level}, \textbf{graph-level} prediction
    (\cref{fig:lec01-tasks}).
\end{itemize}

\begin{figure}[htbp]
  \centering
  \begin{tikzpicture}[q/.style={font=\small\bfseries}]
    \node[vertex] (A) at (0,0) {A};
    \node[target, label={[font=\small, text=ForestGreen!60!black]right:\textbf{?} node level}]
      (B) at (1.7,0.6) {B};
    \node[vertex] (C) at (1.35,-1.3) {C};
    \node[vertex] (D) at (2.5,-0.6) {D};
    \node[vertex] (E) at (3.7,-0.5) {E};
    \node[vertex] (F) at (5.2,0.2) {F};
    \node[vertex] (G) at (5.15,-1.05) {G};
    \node[vertex] (H) at (3.6,-2.1) {H};
    \draw[edge] (A)--(C) (B)--(C) (B)--(D) (C)--(D) (D)--(E) (D)--(H) (E)--(F) (E)--(G);
    \draw[thick, dotted, blue] (A) -- node[q, above left, text=blue, inner sep=1pt]{link level ?} (B);
    \draw[thick, dotted, blue] (C) -- node[q, pos=0.78, below, text=blue]{?} (G);
    \draw[thick, dotted, BurntOrange] (2.4,-0.55) ellipse (4.1 and 2.3);
    \node[font=\small, text=BurntOrange, anchor=west] at (6.2,1.2) {graph level ?};
  \end{tikzpicture}
  \caption{Prediction levels on the example graph (slides 42, 52): label of node B; is
  there a link A--B (or C--G)?; label of the whole graph.}\label{fig:lec01-tasks}
\end{figure}

% =====================================================================
\section{Node-level tasks}\label{sec:lec01-node}

\subsection{Node classification and node-level structure}\label{sec:lec01-node-class}
\slides{44--45}

\begin{itemize}
  \item \term{Node classification}: some nodes labelled (two classes), others unknown
    (``?''); ML assigns labels to all nodes.
  \item \textbf{Node-level network structure. Goal:} characterize the structure and
    position of a node in the network:
    \begin{itemize}
      \item node degree;
      \item node importance \& position -- e.g.\ number of shortest paths passing through a
        node; avg.\ shortest path length to other nodes;
      \item substructures around the node.
    \end{itemize}
    (Slide graph = \cref{fig:lec01-tasks} with A--B as an edge.)
\end{itemize}

\subsection{Example (1): user churn}\label{sec:lec01-churn}
\slides{46}

\begin{itemize}
  \item Users--Sales--Products graph of \cref{fig:lec01-rdl}; predict per user
    \textbf{Churn} vs.\ \textbf{Active}.
  \item \textbf{Training labels together with timestamps are attached to the graph}:
    $u_1$, $u_2$: Churn (May '23), Active (June '23); $u_3$: Active (June '23); $u_4$:
    Active (April '23), Churn (May '23), Active (June '23).
\end{itemize}

\subsection{Example (2): protein folding}\label{sec:lec01-alphafold}
\slides{47--50}

\begin{itemize}
  \item \textbf{Task:} computationally predict a protein's \textbf{3D structure} based solely
    on its amino acid \textbf{sequence}: \textbf{for each node predict its 3D coordinates}.
  \item Experimental result vs.\ computational prediction: T1037\,/\,6vr4 (RNA polymerase
    domain) \textbf{90.7 GDT}; T1049\,/\,6y4f (adhesin tip) \textbf{93.3 GDT}.
  \item \textbf{AlphaFold impact:} median free-modelling accuracy (GDT) at CASP: CASP7--12
    (2006--2016) $\approx$30--40; AlphaFold (CASP13, 2018) $\approx$58; AlphaFold~2 (CASP14,
    2020) $\approx$87 (read off the bar chart). Headlines: ``close to solving one of
    biology's greatest challenges'', ``could turbocharge drug discovery'', etc.
    \textbf{Nobel Prize in Chemistry 2024} (Oct.\ 9, 2024): David Baker (Univ.\ of
    Washington; computational protein design), Demis Hassabis and John M.\ Jumper (Google
    DeepMind; protein structure prediction) -- the second Nobel that week to involve AI.
  \item \textbf{Key idea: ``spatial graph''.} \textbf{Nodes:} amino acids in a protein
    sequence. \textbf{Edges:} proximity between amino acids (residues).
  \item Architecture: protein sequence $\to$ genetics search \& embed $\to$ MSA embedding
    (\emph{sequence--residue edges}); embed \& outer sum $\to$ \emph{residue--residue edges}
    (= spatial graph; gives pairwise distances); both $\to$ structure module $\to$ 3D
    structure + confidence score.
\end{itemize}

% =====================================================================
\section{Link-level tasks}\label{sec:lec01-link}

\subsection{Link-level prediction task}\label{sec:lec01-linkpred}
\slides{52}

\begin{itemize}
  \item Predict \textbf{new/missing/unknown links} based on the existing links.
  \item At test time, node pairs (with no existing links) are \textbf{ranked}, and the
    \textbf{top $K$} node pairs are predicted.
  \item Task: make a prediction for \textbf{a pair of nodes} (A--B, C--G in
    \cref{fig:lec01-tasks}).
\end{itemize}

\subsection{Example (1): recommender systems}\label{sec:lec01-recsys}
\slides{53--54}

\begin{itemize}
  \item Users interact with items (watch movies, buy merchandise, listen to music).
    \textbf{Nodes:} users and items. \textbf{Edges:} user--item interactions.
    \textbf{Goal:} recommend items users might like (``you might also like'' links).
  \item \textbf{PinSage} (Ying et al., KDD 2018): \textbf{task:} recommend related pins to
    users (query pin: cake $\to$ another cake = successful, sweater = bad recommendation).
    Learn node embeddings $z_i$ such that
    \begin{equation}\label{eq:lec01-pinsage}
      d(z_{\text{cake1}}, z_{\text{cake2}}) < d(z_{\text{cake1}}, z_{\text{sweater}})
    \end{equation}
    $\Rightarrow$ \textbf{predict whether two nodes in a graph are related}. (Figure: pins
    linked in a graph; a GNN computation graph produces the embedding $z$ of a pin.)
\end{itemize}

\subsection{Example (2): drug side effects}\label{sec:lec01-polypharmacy}
\slides{55--57}

\begin{itemize}
  \item Many patients take multiple drugs to treat complex or co-existing diseases:
    \textbf{46\% of people ages 70--79 take more than 5 drugs}; many patients take more than
    20 drugs (heart disease, depression, insomnia, etc.).
  \item \textbf{Task:} given a pair of drugs predict adverse side effects (e.g.\ one side
    effect with 30\% prob., another with 65\% prob.).
  \item \textbf{Biomedical graph link prediction} (Zitnik et al., Bioinformatics 2018):
    \textbf{nodes:} drugs \& proteins; \textbf{edges:} interactions -- drug--protein,
    protein--protein, and drug--drug side-effect edges ($r_1$ = gastrointestinal bleed,
    $r_2$ = bradycardia).
  \item \textbf{Query:} how likely will Simvastatin and Ciprofloxacin, when taken together,
    break down muscle tissue?
\end{itemize}

\begin{caution}
  Slide 56 defines $r_2$ = bradycardia in the legend, but labels the queried
  Simvastatin--Ciprofloxacin edge ``$r_2$ (breakdown of muscle tissue)'': the query is a
  different side effect.
\end{caution}

\textbf{Results: de novo predictions} (top 10; 5 with literature evidence):
\begin{center}\small
\begin{tabular}{@{}r l l l l@{}}
  \toprule
  Rank & Drug $c$ & Drug $d$ & Side effect $r$ & Evidence found \\
  \midrule
  1 & Pyrimethamine & Aliskiren & Sarcoma & Stage et al.\ 2015 \\
  2 & Tigecycline & Bimatoprost & Autonomic neuropathy & \\
  3 & Omeprazole & Dacarbazine & Telangiectases & \\
  4 & Tolcapone & Pyrimethamine & Breast disorder & Bicker et al.\ 2017 \\
  5 & Minoxidil & Paricalcitol & Cluster headache & \\
  6 & Omeprazole & Amoxicillin & Renal tubular acidosis & Russo et al.\ 2016 \\
  7 & Anagrelide & Azelaic acid & Cerebral thrombosis & \\
  8 & Atorvastatin & Amlodipine & Muscle inflammation & Banakh et al.\ 2017 \\
  9 & Aliskiren & Tioconazole & Breast inflammation & Parving et al.\ 2012 \\
  10 & Estradiol & Nadolol & Endometriosis & \\
  \bottomrule
\end{tabular}
\end{center}
\begin{itemize}
  \item Case report supporting rank 8: ``Severe Rhabdomyolysis due to Presumed Drug
    Interactions between Atorvastatin with Amlodipine and Ticagrelor''.
\end{itemize}

% =====================================================================
\section{Graph-level tasks}\label{sec:lec01-graph}

\subsection{Graph-level prediction}\label{sec:lec01-graphpred}
\slides{59}

\begin{itemize}
  \item \textbf{Goal:} make a prediction for an \textbf{entire graph} or a \textbf{subgraph}
    of the graph. (Slide example: two clusters $\{A,B,C,D\}$ and $\{E,F,G,H\}$ joined by
    D--E.)
\end{itemize}

\subsection{Example (1): traffic prediction}\label{sec:lec01-traffic}
\slides{60--62}

\begin{itemize}
  \item Google Maps, Stanford $\to$ UC Berkeley: via I-880\,N 51 min (38.9 mi, fastest);
    via I-280\,N 52 min (46.2 mi); via CA-84\,E and I-880\,N 52 min (41.1 mi).
  \item \textbf{Road network as a graph:} \textbf{nodes:} road segments (each with a
    feature vector); \textbf{edges:} connectivity between road segments;
    \textbf{prediction:} time of arrival (ETA) from A to B.
  \item \textbf{Predicting ETA with GNNs} -- \textbf{used in Google Maps}: anonymised travel
    data $\xrightarrow{\text{analysed}}$ supersegments $\xrightarrow{\text{training data}}$
    GNN $\xrightarrow{\text{predictions}}$ routes ranked by ETA (candidate user routes
    A$\to$B from the Google Maps routing system) $\xrightarrow{\text{surfaced}}$ Google
    Maps API / app.
\end{itemize}

\subsection{Example (2): drug discovery}\label{sec:lec01-antibiotics}
\slides{63--64}

\begin{itemize}
  \item \textbf{Antibiotics are small molecular graphs.} \textbf{Nodes:} atoms.
    \textbf{Edges:} chemical bonds. (Slide: $\beta$-lactam families -- penicillins,
    cephalosporins, cephamycins, oxacephems, clavulanic acid, penems, carbapenems,
    nocardicin, monobactams.)
  \item \textbf{Deep learning for antibiotic discovery} (Stokes et al., Cell 2020): a GNN
    \textbf{graph classification model}; predict promising molecules from a pool of
    candidates.
    \begin{itemize}
      \item Model: directed message passing neural network (new vector of bond $2\!\to\!1$
        from concat of bond $2\!\to\!1$ and the sum over bonds into atom 2: $3\!\to\!2$,
        $4\!\to\!2$).
      \item Training set ($10^4$ molecules) $\to$ ML $\to$ predictions \& model validation
        $\to$ iterative re-training; large-scale predictions (upper limit $10^8$+) over the
        chemical landscape $\to$ lead identification \& optimization.
      \item vs.\ conventional small-molecule screening: upper limit $10^5$--$10^6$, hit
        validation 1--3\% hit rate.
    \end{itemize}
\end{itemize}

\subsection{Example (3): physics simulation}\label{sec:lec01-physics}
\slides{65--66}

\begin{itemize}
  \item Sanchez-Gonzalez et al., ICML 2020. \textbf{Physical simulation as a graph:}
    \textbf{nodes:} particles; \textbf{edges:} interaction between particles.
  \item \textbf{Graph evolution task. Goal:} predict how a graph will evolve over time.
    A learned simulator $s_\theta$ (learned $d_\theta$ + Update) is applied repeatedly:
\end{itemize}
\begin{equation}\label{eq:lec01-rollout}
  X^{t_0} \xrightarrow{\;s_\theta\;} \cdots \xrightarrow{\;s_\theta\;} \tilde X^{t_K}
\end{equation}
\begin{itemize}
  \item $d_\theta$ = encoder--processor--decoder (\cref{fig:lec01-epd}):
    (c) \textbf{construct graph}: $\vx_i \to$ node $\vv_i^0$, edges $\ve_{i,j}^0$;
    (d) \textbf{pass messages}: $(\vv_i^m, \ve_{i,j}^m) \to (\vv_i^{m+1}, \ve_{i,j}^{m+1})$;
    (e) \textbf{extract dynamics info}: $\vv_i^M \to \vy_i$.
\end{itemize}

\begin{figure}[htbp]
  \centering
  \begin{tikzpicture}[op/.style={circle, draw, thick, inner sep=0pt, minimum size=4.5mm}]
    \node (X) at (0,0) {$X$};
    \node[block] (enc) at (1.45,0) {Encoder};
    \node (G0) at (2.95,0) {$G^0$};
    \node[op] (p1) at (4.55,0) {$+$};
    \node[block] (gn1) at (3.75,0.95) {GN$^1$};
    \node (G1) at (5.5,0) {$G^1$};
    \node (dots) at (6.3,0) {$\cdots$};
    \node (GM1) at (7.2,0) {$G^{M-1}$};
    \node[op] (pM) at (8.95,0) {$+$};
    \node[block] (gnM) at (8.15,0.95) {GN$^M$};
    \node (GM) at (10.0,0) {$G^M$};
    \node[block] (dec) at (11.5,0) {Decoder};
    \node (Y) at (12.95,0) {$Y$};
    \draw[arrow] (X)--(enc); \draw[arrow] (enc)--(G0); \draw[arrow] (G0)--(p1);
    \draw[arrow] (p1)--(G1);
    \draw[arrow] (G0.north) |- (gn1.west); \draw[arrow] (gn1.east) -| (p1.north);
    \draw[arrow] (GM1)--(pM); \draw[arrow] (pM)--(GM);
    \draw[arrow] (GM1.north) |- (gnM.west); \draw[arrow] (gnM.east) -| (pM.north);
    \draw[arrow] (GM)--(dec); \draw[arrow] (dec)--(Y);
    \begin{scope}[on background layer]
      \fill[black!6, rounded corners=3pt] (3.2,-0.45) rectangle (9.45,1.45);
    \end{scope}
    \node[font=\scriptsize\scshape] at (6.3,1.2) {Processor};
  \end{tikzpicture}
  \caption{Learned simulator $d_\theta$ (slide 66b): encoder, $M$ message-passing blocks
  GN$^m$ with residual $\oplus$, decoder.}\label{fig:lec01-epd}
\end{figure}

\subsection{Application: weather forecasting}\label{sec:lec01-weather}
\slides{67}

\begin{itemize}
  \item \textbf{GraphCast:} (a) input weather state $X^{\le t}$ $\to$ (b) GraphCast predicts
    the next state $\hat X^{t+1}$ $\to$ (c) rolling out a forecast GC $\to\cdots\to$ GC
    gives $\hat X^{t+T}$.
  \item Inside: (d) encoder (grid $\to$ mesh), (e) processor, (f) decoder (mesh $\to$ grid);
    (g) \textbf{simultaneous multi-mesh message passing} over meshes $M^0,\dots,M^6$
    (refined icosahedra).
  \item Linked article: GraphCast outperforms the world's best medium-range weather
    forecasting.
\end{itemize}

% =====================================================================
\par\vspace{0pt plus 6\baselineskip}\penalty-300\vspace{0pt plus -6\baselineskip}
\section{Summary}\label{sec:lec01-summary}
\slides{68}

\textbf{ML in the language of graphs} (on the Users--Sales--Products graph,
\cref{fig:lec01-rdl}):
\begin{itemize}
  \item \textbf{Node-level:} churn, life-time value, next best action.
  \item \textbf{Link-level:} product affinity, recommendations.
  \item \textbf{Graph-level:} fraud, money laundering.
\end{itemize}

% ---------------------------------------------------------------------
\bigskip\phantomsection
\addcontentsline{toc}{section}{Lecture 1 at a glance}
\begin{summarybox}[Lecture 1 at a glance]
  \begin{itemize}
    \item Graphs = general language for entities with relations; databases are graphs
      (RDL: rows = nodes, key links = edges). Explicitly modeling relationships
      $\Rightarrow$ better performance.
    \item DL toolbox is built for sequences \& grids; networks have arbitrary size/topology,
      no fixed node ordering, are dynamic and multimodal $\Rightarrow$ graphs are the new
      frontier of DL.
    \item Representation: choose nodes and edges -- the choice of links fixes the questions
      you can study. Heterogeneous graph $G=(V,E,R,T)$, edges $(v_i,r,v_j)$; directed vs.\
      undirected; weights, properties, types, attributes; bipartite graphs and folded
      networks.
    \item Tasks: node level (classification; churn; AlphaFold 3D coordinates), link level
      (rank non-linked pairs, top $K$; PinSage
      $d(z_{\text{cake1}},z_{\text{cake2}})<d(z_{\text{cake1}},z_{\text{sweater}})$;
      polypharmacy side effects), graph/subgraph level (ETA in Google Maps; antibiotic
      graph classification; physics simulation as graph evolution; GraphCast weather forecasting).
  \end{itemize}
\end{summarybox}
